\documentclass[12pt,a4paper]{article}

\usepackage[utf8]{inputenc}
\usepackage[T1]{fontenc}
\usepackage[brazil]{babel}
\usepackage{geometry}
\geometry{top=3cm,bottom=2cm,left=3cm,right=2cm}
\usepackage{setspace}
\onehalfspacing
\usepackage{microtype}
\usepackage{amsmath,amssymb}
\usepackage{graphicx}
\usepackage{booktabs}
\usepackage{tabularx}
\usepackage{xcolor}
\usepackage{listings}
\usepackage{hyperref}

\hypersetup{
    colorlinks=true,
    linkcolor=blue!70!black,
    citecolor=blue!70!black,
    urlcolor=blue!70!black
}

\lstset{
    basicstyle=\ttfamily\small,
    keywordstyle=\color{blue!80!black}\bfseries,
    stringstyle=\color{teal!80!black},
    commentstyle=\color{gray}\itshape,
    numbers=none,
    frame=single,
    breaklines=true,
    showstringspaces=false,
    tabsize=2
}

\title{
    \textbf{Atividade 3: Comunicação em Sistemas Distribuídos}\\
    \large Um Estudo Comparativo de Chat Distribuído com Sockets TCP, RPC (gRPC) e Java RMI
}
\author{
    \textbf{Pedro Bazza} \and \textbf{Matheus} \\
    \textit{Disciplina de Sistemas Distribuídos}
}
\date{Outubro de 2026}

\begin{document}

\maketitle

\begin{abstract}
Este trabalho apresenta o projeto, a implementação e a análise comparativa de uma aplicação distribuída de chat cliente-servidor desenvolvida em Java (JDK 26) sob três paradigmas distintos de comunicação em rede: \textbf{Sockets TCP}, \textbf{Remote Procedure Call (RPC via gRPC)} e \textbf{Remote Method Invocation (Java RMI)}. Para garantir máxima coesão e isolamento tecnológico, a aplicação foi estruturada seguindo os princípios da Arquitetura Hexagonal (\textit{Ports \& Adapters}) em conjunto com o padrão de projeto \textit{Strategy}, permitindo a substituição transparente do mecanismo de transporte em tempo de execução via argumentos de linha de comando. Adicionalmente, toda a infraestrutura foi containerizada através do Docker e Docker Compose, assegurando reprodutibilidade e portabilidade. O relatório detalha as características arquiteturais de cada versão, avalia o nível de abstração, desempenho e complexidade, e discute as dificuldades práticas superadas durante o desenvolvimento.
\end{abstract}

\newpage
\tableofcontents
\newpage

\section{Introdução e Definição da Aplicação}

O desenvolvimento de sistemas distribuídos contemporâneos impõe desafios contínuos relativos ao desacoplamento entre módulos de software e à escolha do modelo de comunicação adequado para os requisitos de latência, escalabilidade, interoperabilidade e manutenibilidade. 

Com o intuito de explorar pragmaticamente tais aspectos, este projeto propôs a implementação de um **Sistema de Chat Distribuído Cliente-Servidor Multi-Dispositivo**, no qual múltiplos clientes podem se conectar simultaneamente a um servidor central, enviar mensagens de texto e recebê-las de forma assíncrona e contínua via \textit{broadcast}.

\subsection{Justificativa da Aplicação Escolhida}
A aplicação de chat foi selecionada por demandar simultaneamente características fundamentais de sistemas distribuídos:
\begin{enumerate}
    \item \textbf{Bidirecionalidade Contínua:} O cliente não apenas requisita dados, mas deve ser capaz de receber mensagens assíncronas do servidor no momento exato em que são publicadas por outros clientes.
    \item \textbf{Concorrência e Compartilhamento de Estado:} O servidor deve gerenciar múltiplos canais ativos concorrentemente na memória RAM sem corromper o estado dos clientes conectados.
    \item \textbf{Heterogeneidade Conceitual:} O problema de transmitir uma mensagem para múltiplos participantes evidencia claramente como cada tecnologia lida com fluxo de bytes (Sockets), chamadas remotas de procedimentos (RPC) e referências a objetos remotos (RMI).
\end{enumerate}

---

\section{Arquitetura Geral do Sistema}

Para viabilizar que as três tecnologias coexistissem no mesmo projeto sem duplicação de lógica de apresentação e domínio, adotou-se a **Arquitetura Hexagonal (\textit{Ports \& Adapters})** combinada com o padrão de projeto **\textit{Strategy}**.

\subsection{Estrutura de Camadas}
\begin{itemize}
    \item \textbf{Core de Domínio (\textit{Domain}):} Contém as entidades nucleares do sistema, livres de qualquer dependência externa ou tecnológica:
    \begin{itemize}
        \item \texttt{Message}: Representa a mensagem de texto contendo remetente (\textit{sender}), conteúdo textual (\textit{content}), carimbo de data/hora (\textit{timestamp}) e destinatário (\textit{recipient}). Implementa \texttt{Serializable} para atender à serialização nativa do Java.
        \item \texttt{User}: Encapsula o identificador ou nome de usuário (\textit{username}) do participante conectado.
    \end{itemize}
    \item \textbf{Portas (\textit{Ports}):} Definem os contratos abstratos de comunicação:
    \begin{itemize}
        \item \texttt{ChatServerPort}: Contrato do servidor, exigindo operações de ciclo de vida (\texttt{start}, \texttt{stop}) e disseminação (\texttt{broadcast}).
        \item \texttt{ChatClientPort}: Contrato do cliente, exigindo operações de conexão (\texttt{connect}), envio (\texttt{sendMessage}) e encerramento (\texttt{disconnect}).
        \item \texttt{MessageListener}: Interface funcional atuando como \textit{Observer}, pela qual as camadas de rede notificam a interface de usuário sempre que uma nova mensagem for entregue.
    \end{itemize}
    \item \textbf{Adaptadores de Rede (\textit{Adapters}):} Implementações concretas das portas para cada protocolo: \texttt{sockets}, \texttt{grpc} e \texttt{rmi}.
    \item \textbf{Fábrica de Rede (\textit{NetworkFactory}):} Ponto centralizado de injeção que recebe a \textit{string} correspondente ao protocolo via CLI (\texttt{"socket"}, \texttt{"grpc"} ou \texttt{"rmi"}) e instancia o adaptador correto.
    \item \textbf{Interface de Linha de Comando (\textit{CLI UI}):} Implementa \texttt{MessageListener} para exibir mensagens na saída padrão sem bloquear o laço de leitura de entrada do usuário via terminal.
\end{itemize}

\subsection{Containerização com Docker}
Conforme decisão técnica do projeto, toda a aplicação foi containerizada empregando o **JDK 26** em um \textit{build multi-stage}:
\begin{itemize}
    \item No estágio de compilação (\textit{builder}), o Maven gera os stubs a partir da definição do Protocol Buffers (\texttt{.proto}) e empacota todas as dependências em um \textit{shaded JAR} executável.
    \item No estágio final de execução (\textit{runtime}), uma imagem enxuta baseada em \texttt{eclipse-temurin:26} expõe as portas de serviço e fornece os \textit{entrypoints} unificados para o servidor (\texttt{ServerApp}) e clientes (\texttt{ClientApp}).
    \item O \texttt{docker-compose.yml} orquestra instâncias independentes dos três servidores e facilita a execução dos clientes no modo interativo.
\end{itemize}

---

\section{Arquitetura e Implementação das Três Versões}

\subsection{Versão 1: Sockets TCP}
A versão baseada em sockets emprega a API nativa \texttt{java.net.ServerSocket} e \texttt{java.net.Socket}, operando na camada de transporte sobre o protocolo TCP (\textit{Transmission Control Protocol}), garantindo entrega confiável, ordenada e com verificação de erros.

\subsubsection{Mecanismo de Comunicação}
\begin{itemize}
    \item \textbf{Servidor (\texttt{SocketServerAdapter}):} Escuta requisições na porta designada (padrão \texttt{8080}). Ao aceitar uma conexão, delega o canal para uma thread dedicada executada por um \texttt{CachedThreadPool} gerenciando instâncias de \texttt{SocketClientHandler}.
    \item \textbf{Serialização:} O transporte de objetos \texttt{Message} entre cliente e servidor foi realizado via \texttt{ObjectOutputStream} e \texttt{ObjectInputStream}. Após o envio de cada objeto, executa-se \texttt{flush()} e \texttt{reset()} para descarregar o buffer do socket e limpar o cache de instâncias do fluxo de serialização.
    \item \textbf{Recepção Assíncrona no Cliente (\texttt{SocketClientAdapter}):} Ao conectar, o cliente inicia uma thread em segundo plano (\textit{daemon}) dedicada exclusivamente a ler o \texttt{ObjectInputStream} e acionar o callback \texttt{MessageListener.onMessageReceived()}.
\end{itemize}

\subsection{Versão 2: RPC com gRPC e Protocol Buffers}
O paradigma RPC (\textit{Remote Procedure Call}) abstrai a rede de modo que invocar um método em um servidor remoto seja sintaticamente idêntico a chamar uma função local. Para esta versão, utilizou-se o framework de alta performance **gRPC** sobre o protocolo **HTTP/2**, com especificação formal via **Protocol Buffers v3**.

\subsubsection{Definição do Contrato (\texttt{chat.proto})}
O contrato de comunicação foi definido formalmente:
\begin{lstlisting}[language=protobuf]
syntax = "proto3";
package com.pedrobazza.distsyschat.adapters.grpc;

service ChatService {
    rpc StreamMessages(ConnectRequest) returns (stream ChatMessage);
    rpc SendMessage(ChatMessage) returns (SendResponse);
    rpc Disconnect(DisconnectRequest) returns (DisconnectResponse);
}

message ConnectRequest { string username = 1; }
message ChatMessage {
    string sender = 1;
    string content = 2;
    int64 timestamp = 3;
    string recipient = 4;
}
message SendResponse { bool success = 1; string message = 2; }
message DisconnectRequest { string username = 1; }
message DisconnectResponse { bool success = 1; }
\end{lstlisting}

\subsubsection{Mecanismo de Comunicação}
\begin{itemize}
    \item \textbf{Streaming Unidirecional Servidor-Cliente (\textit{Server-Streaming RPC}):} O método remoto \texttt{StreamMessages} permite que o cliente envie seu nome de usuário e receba um fluxo contínuo de eventos (\texttt{stream ChatMessage}) via \texttt{StreamObserver}. O servidor mantém as instâncias ativas de \texttt{StreamObserver} em memória para enviar broadcasts instantâneos.
    \item \textbf{Invocação Unária (\textit{Unary RPC}):} Os métodos \texttt{SendMessage} e \texttt{Disconnect} realizam chamadas diretas com retorno estruturado, manipuladas pelo stub bloqueante (\texttt{ChatServiceBlockingStub}).
\end{itemize}

\subsection{Versão 3: Java RMI (Remote Method Invocation)}
O Java RMI é a tecnologia proprietária da plataforma Java para comunicação baseada em **Objetos Distribuídos**. Diferente de Sockets e RPC convencionais, o RMI permite invocar métodos em instâncias de objetos que residem em uma JVM remota, com transmissão transparente de parâmetros e exceções por meio de stubs e skeletons gerados dinamicamente.

\subsubsection{Mecanismo de Comunicação}
\begin{itemize}
    \item \textbf{Interface Remota (\texttt{RmiRemoteService}):} Estende \texttt{java.rmi.Remote} e declara os métodos remotos que o servidor disponibiliza:
    \begin{lstlisting}[language=Java]
public interface RmiRemoteService extends Remote {
    interface ClientCallback extends Remote {
        void receiveMessage(Message message) throws RemoteException;
    }
    void registerClient(String username, ClientCallback callback) throws RemoteException;
    void unregisterClient(String username) throws RemoteException;
    void postMessage(Message message) throws RemoteException;
}
    \end{lstlisting}
    \item \textbf{Registro de Nomes (\textit{RMI Registry}):} O servidor instancia um serviço remoto (\texttt{RmiRemoteServiceImpl}) estendendo \texttt{UnicastRemoteObject} e registra sua referência no \textit{RMI Registry} na porta \texttt{1099} através do nome chave \texttt{"ChatService"}.
    \item \textbf{Callback Remoto para Broadcast:} Como o RMI é classicamente orientado a requisições do cliente para o servidor, a bidirecionalidade do chat exigiu a criação de um **Callback Remoto**: o cliente exporta um objeto remoto anônimo (\texttt{ClientCallback}) e envia essa referência remota ao servidor no momento do registro. Ao receber uma mensagem, o servidor itera sobre a coleção de stubs de callback dos clientes e invoca remotamente o método \texttt{receiveMessage(message)} de cada um.
\end{itemize}

---

\section{Comparação Técnica Entre as Três Abordagens}

A Tabela~\ref{tab:comparacao} resume as principais características comparativas entre Sockets TCP, gRPC e Java RMI no contexto do desenvolvimento da aplicação de chat.

\begin{table}[htbp]
\centering
\small
\caption{Comparação entre Sockets, gRPC e Java RMI.}
\label{tab:comparacao}
\begin{tabularx}{\textwidth}{lXXX}
\toprule
\textbf{Critério} & \textbf{Sockets TCP} & \textbf{gRPC (RPC)} & \textbf{Java RMI} \\
\midrule
\textbf{Nível de Abstração} & Baixo (fluxo de bytes bruto ou objetos serializados). & Alto (procedimentos remotos com contrato formal IDL). & Muito Alto (invocação de métodos em instâncias distribuídas). \\
\midrule
\textbf{Protocolo de Transporte} & TCP puro com portas customizadas. & HTTP/2 com multiplexação de conexões. & TCP / JRMP (\textit{Java Remote Method Protocol}). \\
\midrule
\textbf{Formato dos Dados} & Binário via serialização Java nativa. & Binário compacto via Protocol Buffers. & Binário via serialização nativa de classes Java. \\
\midrule
\textbf{Bidirecionalidade} & Manual: controle de threads separadas de leitura e escrita. & Nativa: suporte embutido a streaming unidirecional e bidirecional. & Indireta: requer exportação de objetos de \textit{callback} pelo cliente. \\
\midrule
\textbf{Interoperabilidade} & Média (depende da especificação manual do protocolo). & Excelente (agnóstico a plataformas e linguagens). & Nula (restrito e acoplado ao ecossistema JVM/Java). \\
\midrule
\textbf{Complexidade de Código} & Alta para sincronização e tratamento de exceções de rede. & Média (requer etapa de compilação de stubs, mas simplifica o código). & Baixa para chamadas simples; Média/Alta para gerenciar callbacks e rede. \\
\bottomrule
\end{tabularx}
\end{table}

\subsection{Nível de Abstração Oferecido}
\begin{itemize}
    \item \textbf{Sockets:} Exige que o desenvolvedor gerencie explicitamente todos os aspectos da comunicação: abertura e fechamento de sockets, instanciação de buffers de entrada/saída, tratamento de conexões prematuramente encerradas e criação de laços com multithreading.
    \item \textbf{gRPC:} Fornece abstração focada no negócio. O desenvolvedor interage com classes fortemente tipadas geradas pelo compilador, não precisando lidar com sockets TCP ou detalhes de transporte HTTP/2.
    \item \textbf{RMI:} Oferece a mais alta transparência sintática para desenvolvedores Java, pois a invocação remota parece uma chamada de método Java comum (\texttt{service.postMessage(msg)}), encapsulando a complexidade da rede nos stubs.
\end{itemize}

\subsection{Facilidade de Uso e Complexidade}
Embora Sockets não exijam bibliotecas externas, a complexidade de manter estados concorrentes e gerenciar threads de leitura/escrita eleva o risco de vazamentos de recursos ou travamentos por concorrência. 

O gRPC possui uma curva de aprendizado inicial vinculada à sintaxe do Protobuf e à configuração dos plugins Maven de geração de código. No entanto, uma vez configurado, o desenvolvimento torna-se extremamente previsível e robusto.

O RMI apresenta a menor complexidade de código inicial para comunicação básica em Java, mas expõe dificuldades relevantes em ambientes modernos com contêineres e NAT, uma vez que requer configuração precisa de hostnames públicos e portas dinâmicas para os stubs exportados.

\subsection{Casos de Uso Práticos}
\begin{itemize}
    \item \textbf{Sockets TCP:} Indicado para jogos online em tempo real, sistemas embarcados com hardware limitado, protocolos proprietários de telecomunicações e serviços em que o controle sobre cada byte transmitido é imperativo.
    \item \textbf{gRPC:} Padrão da indústria para arquiteturas de microsserviços em nuvem, comunicação entre serviços de backend distribuídos em diferentes linguagens (ex: Java, Go, Python, C++) e aplicações móveis que demandam streaming eficiente e baixo consumo de banda.
    \item \textbf{Java RMI:} Adequado para sistemas legados corporativos puramente Java (JEE/EJB) operando dentro de redes locais corporativas fechadas, onde a homogeneidade de tecnologia é garantida.
\end{itemize}

---

\section{Dificuldades Encontradas e Soluções Adotadas}

Durante o desenvolvimento da Etapa 1, alguns desafios técnicos foram identificados e solucionados:

\begin{enumerate}
    \item \textbf{Conflito de Provedores SPI no Empacotamento gRPC com Maven Shade:}
    Ao gerar o JAR executável contendo todas as dependências (\textit{fat JAR}), o gRPC falhava na resolução de nomes do canal com a mensagem de que nenhum \texttt{NameResolverProvider} havia sido encontrado. 
    \textit{Solução:} Adicionou-se o \texttt{ServicesResourceTransformer} à configuração do \texttt{maven-shade-plugin} no \texttt{pom.xml}, garantindo que todos os arquivos declaratórios em \texttt{META-INF/services/} fossem mesclados ao invés de sobrescritos.
    
    \item \textbf{Roteamento de Callbacks RMI em Rede Docker Bridge:}
    Em contêineres Docker, o RMI Registry registrava o endereço interno efêmero do contêiner. Quando o cliente tentava registrar o callback ou o servidor tentava invocar métodos de volta no cliente, ocorria travamento por tempo limite de conexão (\textit{timeout}).
    \textit{Solução:} Configurou-se a propriedade de sistema \texttt{java.rmi.server.hostname} do Java RMI com suporte à leitura de variáveis de ambiente (\texttt{RMI\_HOST}), permitindo que os contêineres se comuniquem corretamente através do DNS interno da rede do Docker Compose.
    
    \item \textbf{Finalização Graciosa de Processos Multithread:}
    O cliente RMI mantinha threads ativas mesmo após o encerramento do chat, impedindo a finalização do processo do terminal.
    \textit{Solução:} Implementou-se o desregistro explícito com \texttt{unexportObject(callback, true)} associado à chamada deliberada de encerramento de threads residuais no fechamento da interface de linha de comando.
\end{enumerate}

---

\section{Conclusão}

A execução prática deste projeto atingiu plenamente os objetivos propostos na Atividade 3. A implementação das três versões de chat com Sockets TCP, gRPC e Java RMI sobre a mesma aplicação evidenciou claramente a evolução dos modelos de comunicação distribuída: desde a manipulação direta de fluxos de bytes até a abstração sofisticada de procedimentos e objetos remotos.

A adoção da Arquitetura Hexagonal com o padrão Strategy demonstrou-se fundamental para desacoplar a regra de negócio das especificidades de cada protocolo, permitindo que os mesmos contratos fossem testados sob diferentes transportes com zero impacto na camada de interface com o usuário. Finalmente, a containerização com Docker assegurou a consistência do ambiente de desenvolvimento no moderno JDK 26, garantindo a rápida validação das três abordagens de sistemas distribuídos.

\end{document}
